\phantomsection
\label{fig:late-han-regions}
\begin{center}
\begin{tikzpicture}[x=.88cm,y=.88cm,font=\small]
  \fill[paper] (0,0) rectangle (12.3,6.65);

  \draw[source!70,line width=.8pt] (.55,2.1) .. controls (3.5,1.72) and
    (7.1,2.55) .. (11.8,1.82);
  \node[text=source,below,font=\scriptsize] at (8.4,1.92) {长江方向（示意）};
  \draw[source!55,dashed] (.75,4.06) .. controls (3.7,4.55) and
    (7.5,3.55) .. (11.7,4.12);
  \node[text=source,above,font=\scriptsize] at (8.7,3.88) {黄河方向（示意）};

  \fill[dynasty!44] (6.08,3.8) ellipse (1.28 and .86);
  \node[align=center,text=white] at (6.08,3.8) {洛阳\\东汉中枢};

  \fill[timeline!26] (2.0,4.78) ellipse (1.36 and .84);
  \node[align=center] at (2.0,4.78) {冀州、广宗\\黄巾主战场};

  \fill[source!24] (3.0,1.05) ellipse (1.35 and .74);
  \node[align=center] at (3.0,1.05) {凉州\\军民与羌胡起事};

  \fill[timeline!27] (8.8,1.16) ellipse (1.55 and .8);
  \node[align=center] at (8.8,1.16) {荆、扬与江淮\\地方军政网络};

  \fill[dynasty!25] (10.46,4.86) ellipse (1.22 and .84);
  \node[align=center] at (10.46,4.86) {冀州、兖州\\袁绍、曹操等};

  \fill[source!28] (5.78,5.48) ellipse (1.18 and .61);
  \node[align=center,font=\scriptsize] at (5.78,5.48) {许县\\献帝与曹操};

  \draw[->,very thick,color=dynasty]
    (5.03,4.16) .. controls (4.3,4.54) and (3.7,4.74) .. (3.31,4.78);
  \node[text=dynasty,above,font=\scriptsize,align=center] at (4.12,5.15)
    {184年调兵\\与地方响应};

  \draw[->,thick,dashed,color=timeline]
    (5.0,3.3) .. controls (4.24,2.38) and (3.55,1.62) .. (3.22,1.53);
  \node[text=timeline,right,font=\scriptsize,align=left] at (3.82,2.4)
    {边郡军费、迁徙\\与地方动员};

  \draw[->,very thick,color=dynasty]
    (6.78,4.39) .. controls (7.35,4.85) and (8.17,5.05) .. (9.17,4.96);
  \node[text=dynasty,above,font=\scriptsize] at (7.88,5.38)
    {州郡与军阀竞争};

  \draw[->,thick,dashed,color=source]
    (6.5,3.1) .. controls (7.14,2.46) and (7.82,1.83) .. (8.12,1.63);
  \node[text=source,right,font=\scriptsize,align=left] at (7.04,2.55)
    {长江、淮河\\限制调兵};

  \draw[->,thick,color=timeline]
    (6.02,4.68) .. controls (5.96,4.94) and (5.89,5.05) .. (5.84,5.08);
  \node[text=timeline,left,font=\scriptsize,align=right] at (5.35,5.03)
    {196年\\迎帝至许};

  \node[draw=source!75,fill=paper,rounded corners=2pt,align=center,
    inner sep=3pt,font=\scriptsize] at (8.86,6.02)
    {黄巾、凉州与继承危机使军队、粮道和皇帝名义\\
    分别被不同地方集团重新接合};
\end{tikzpicture}

\smallskip
\small\textit{图\quad 东汉晚期的几处动员中心与军政网络，约184--200年：据《后汉书·灵帝纪》《献帝纪》、〈皇甫嵩传〉〈董卓传〉，《三国志·武帝纪》及《资治通鉴》卷五十八至六十三，概括黄巾起事、凉州动乱、关东诸军与曹操迎献帝后形成的关系。节点和河流只示相对位置；色块不是州界或实际控制范围，箭线也不重建军队、流民或献帝行幸的精确路线。\quad\hyperlink{evt:EH-0184-YELLOW-TURBAN}{回到“黄巾起事”}}
\end{center}
